% 2011 年天津卷文科第 17 题：四棱锥 P-ABCD，底面 ABCD 为平行四边形，
% O 为 AC 中点，PO⊥底面，M 为 PD 中点。
% 来源 PDF：普通高考/2011/2011天津文.pdf，第 2 页（印面 708）第 1 栏。
% 重绘依据：原图为斜二测式透视，D 在左后、A 在左前、B 在右前、C 在右后；
%   P 在 O 正上方；可见棱 PA、PB、PC、PD、AB、BC 画实线，被遮挡的 AD、CD、AC
%   与辅助线 OM 画虚线。
\begin{tikzpicture}[scale=0.95, semithick,
  every node/.style={inner sep=1.5pt}]
  \coordinate (A) at (0.30,0.20);
  \coordinate (B) at (3.30,0.20);
  \coordinate (C) at (4.10,1.35);
  \coordinate (D) at (1.10,1.35);
  \coordinate (O) at (2.20,0.78);
  \coordinate (P) at (2.20,3.70);
  \coordinate (M) at ($(P)!0.5!(D)$);
  % 底面被遮挡的边与对角线
  \draw[dashed] (A) -- (D) -- (C);
  \draw[dashed] (A) -- (C);
  % 可见棱
  \draw (A) -- (B) -- (C);
  \draw (P) -- (A);
  \draw (P) -- (B);
  \draw (P) -- (C);
  \draw (P) -- (D);
  % 辅助线 OM（中位线，图中虚线）
  \draw[dashed] (O) -- (M);
  % 标注
  \node[left] at ([xshift=-1pt]A) {$A$};
  \node[below] at (B) {$B$};
  \node[right] at (C) {$C$};
  \node[above left] at (D) {$D$};
  \node[above] at (P) {$P$};
  \node[left] at ([xshift=1pt]M) {$M$};
  \node[below] at (O) {$O$};
\end{tikzpicture}
